\section{Implementazione Java}

\subsection{Struttura del progetto Maven}

Il progetto Maven (\texttt{java/}) gira su Java 17 e usa Spring Boot con Vaadin 24 per
la UI, JUnit 5 e Mockito per i test e JaCoCo per la copertura. Il cuore del programma sta
nel package \texttt{it.tvsw.smartparking.core}, che contiene:

\begin{itemize}
    \item \texttt{StatoSistema}, \texttt{TipoUtente}, \texttt{TipoPosto}: enum identici ai domini ASM;
    \item \texttt{Parcheggio}: nucleo con contratti JML (Sezione~\ref{sec:jml-ref});
    \item \texttt{SensorInput}: raccoglie i 6 input monitorati di uno step;
    \item \texttt{Display}: interfaccia di notifica (usata per Mockito);
    \item \texttt{SmartParkingFSM}: la macchina a stati, un metodo \texttt{step(SensorInput)} che replica la main rule ASM.
\end{itemize}

\subsection{La macchina a stati}

\begin{lstlisting}[language=Java, caption={step() -- dispatch sullo stato corrente}]
public void step(SensorInput input) {
    switch (stato) {
        case IDLE:     gestisciIdle(input); break;
        case CHK_IN:   gestisciChkIn(); break;
        case VER:      gestisciVer(input); break;
        case NEG:      gestisciNeg(input); break;
        case INGR:     gestisciIngr(input); break;
        case CHK_OUT:  gestisciChkOut(input); break;
        case TARIF:    gestisciTarif(input); break;
        case USC:      gestisciUsc(input); break;
    }
}
\end{lstlisting}

Una nota sul design: nel modello ASM il contatore viene decrementato solo al transito
vero e proprio in INGR, mentre qui la classe \texttt{Parcheggio} decide e decrementa
tutto in un colpo solo, dentro \texttt{assegnaPosto} (che viene chiamato da
\texttt{gestisciVer}). Negli scenari di test del progetto questa differenza non si vede,
e l'ho scelta apposta per tenere il codice più semplice; ne riparlo meglio
nell'osservazione dedicata in Sezione~\ref{sec:analisi-statica}.

\subsection{Interfaccia utente (Vaadin)}

Per la UI ho creato la vista \texttt{MainView}, che mostra lo stato corrente e i posti
ancora liberi. I sensori li ho simulati con dei bottoni: ``Arriva auto
STD/DISABILE/ABBONATO'', ``Transito'', ``Pagamento'', ``Auto va via'' e ``Auto in
uscita''. In pratica cliccando i bottoni si fa avanzare la macchina a stati come se
arrivassero gli input dai sensori veri.

\screenshot{vaadin_ui_idle}{UI Vaadin -- stato iniziale}{%
Lanciare \texttt{cd java \&\& mvn spring-boot:run}, aprire
\url{http://localhost:8080}. Screenshot della pagina appena caricata: stato IDLE,
posti standard liberi = 1, posti disabili liberi = 1.}

\screenshot{vaadin_ui_ingresso_std}{UI Vaadin -- dopo un ingresso STD}{%
Cliccare ``Arriva auto STD'' poi ``Transito''. Screenshot che mostra lo stato
aggiornato (IDLE) e posti standard liberi = 0.}

\screenshot{vaadin_ui_neg}{UI Vaadin -- accesso negato a parcheggio pieno}{%
Con il parcheggio pieno (posti standard e disabili a 0), cliccare di nuovo un
bottone ``Arriva auto...''. Screenshot che mostra il messaggio ``Accesso negato''.}

\screenshot{vaadin_ui_tarif}{UI Vaadin -- attesa pagamento}{%
Dopo un ingresso STD e un ``Transito'', cliccare ``Auto in uscita'': lo stato deve
mostrare TARIF (in attesa di pagamento) prima di cliccare ``Pagamento''.}

\subsection{Test end-to-end con Selenium}

La UI l'ho controllata anche con un test end-to-end fatto in Selenium
(\texttt{UISeleniumTest}). L'ho taggato ``ui'' e l'ho escluso dalla build normale, perché
per girare ha bisogno dell'applicazione già avviata e di Chrome. Quello che fa è aprire
la pagina con ChromeDriver in modalità headless, aspettare che Vaadin finisca di
renderizzare lato client (con \texttt{WebDriverWait}), cliccare ``Arriva auto STD'' e
controllare che lo stato mostrato diventi INGR o NEG. Per lanciarlo basta
\texttt{mvn test -Dtest.groups=ui -Dtest.excludedGroups=}, oppure eseguirlo da Eclipse
come un normale test JUnit.

\screenshot{selenium_test_green}{Test Selenium end-to-end superato}{%
Vista JUnit di Eclipse (o terminale Maven) con \texttt{UISeleniumTest} verde,
eseguito con l'applicazione attiva su localhost:8080.}
